\documentclass[../../main.tex]{subfiles}

\begin{document}

\section{The Jordan--Wigner mapping for spinless Fermions}
\label{sec:jordan-wigner}

\subsection{Summary}

The Jordan--Wigner transformation (JWT) can be used to map a fermionic system to a system of
spins. This is essentially done by rewriting
\begin{equation*}
\text{fermion}=\text{string}\times\text{operator},
\end{equation*}
where the ``string'' is some product of operators $\hat Z_i=1-2\hat n_i$ and the ``operator''
is found to obey spin commutation relations.

As references I used
\begin{itemize}
\item \emph{Fermions and Jordan--Wigner String} -- ITensor C\texttt{++} Documentation,
\item \emph{Introduction to Many-Body Physics}, Chapter 4.1 -- Piers Coleman,
\item \emph{The Fermionic canonical commutation relations and the Jordan--Wigner transform}
      -- M.\ Nielsen,
\end{itemize}
which I supplemented with some calculations and my own way of structuring things.

\subsection{Setup: the fermionic Hilbert space}

Our fermionic Hilbert space is defined by (Nielsen):
\begin{enumerate}
\item[i)] operators $c_1,\dots,c_L$ acting on some Hilbert space;
\item[ii)] these operators satisfying the canonical (anti-)commutation relations for
fermions
\begin{equation}
\{c_i^{(+)},c_j^{(+)}\}=0\qquad\&\qquad\{c_i,c_j^\dagger\}=\delta_{ij}\,\mathds{1}.
\end{equation}
\end{enumerate}
Many things directly follow from this, such as a possible basis for the Hilbert space, namely
the occupation number basis
\begin{equation}
\ket{\vec n}=\ket{n_1,n_2,\dots,n_L}\equiv\left(\prod_{i=1}^{L}(c_i^\dagger)^{n_i}\right)\ket{0},
\end{equation}
where $\hat n_i\ket{\vec n}=c_i^\dagger c_i\ket{\vec n}=n_i\ket{\vec n}$ with $n_i=0,1$ and
\begin{equation}
[\hat n_i,\hat n_j]=0.\label{eq:jw-nn}
\end{equation}

\subsection{The transformation}

The JWT introduces new operators $a_j$ by writing
\begin{equation}
c_j\equiv\left(\prod_{i=1}^{j-1}Z_i\right)a_j\qquad\&\qquad
c_j^\dagger\equiv a_j^\dagger\left(\prod_{i=1}^{j-1}Z_i\right), \label{eq:jw-def}
\end{equation}
with $Z_j=(-1)^{n_j}=1-2n_j$ and $n_j=c_j^\dagger c_j$. The ordering of the $Z_i$ does not
matter since $[Z_i,Z_j]=0$ due to \cref{eq:jw-nn}. Inverting \cref{eq:jw-def} using $Z_i^2=\mathds{1}$ and $[Z_i,Z_j]=0$ yields
\begin{equation}
a_j=\left(\prod_{i=1}^{j-1}Z_i\right)c_j\qquad\&\qquad
a_j^\dagger=c_j^\dagger\left(\prod_{i=1}^{j-1}Z_i\right). \label{eq:jw-inv}
\end{equation}

\subsection{Algebra of the new operators}

One can show that these new operators obey
\begin{align}
[a_i^{(+)},a_j^{(+)}]&=0 &&\forall\,i,j,\\
[a_i,a_j^\dagger]&=0 &&\forall\,i\neq j,\\
[a_i,a_i^\dagger]&=1-2n_i &&\text{with }n_i=c_i^\dagger c_i,
\end{align}
using that
\begin{equation}
\{Z_i,c_j^{(+)}\}=0 \label{eq:jw-Zc}
\end{equation}
The corresponding proofs can be found in \cref{subsec:proofs_jwt}.

\subsection{The occupation-number basis in the new operators}

The basis states of the fermionic Hilbert space can be written as
\begin{align}
\ket{n_1,n_2,\dots,n_L}
&\equiv\left(\prod_{i=1}^{L}(c_i^\dagger)^{n_i}\right)\ket{0}
&&\text{w.\ }n_i=0,1\notag\\
&=\left\{\prod_{i=1}^{L}\left[a_i^\dagger\left(\prod_{j=1}^{i-1}Z_j\right)\right]^{n_i}\right\}\ket{0}
&&(c_1^\dagger=a_1^\dagger)\notag\\
&=(a_1^\dagger)^{n_1}(a_2^\dagger Z_1)^{n_2}(a_3^\dagger Z_2 Z_1)^{n_3}\cdots
  (a_L^\dagger Z_{L-1}\cdots Z_1)^{n_L}\ket{0}\notag\\
&=\left(\prod_{i=1}^{L}(a_i^\dagger)^{n_i}\right)
  (Z_1)^{n_2}(Z_2 Z_1)^{n_3}\cdots(Z_{L-1}\cdots Z_1)^{n_L}\ket{0}\notag\\
&=\left(\prod_{i=1}^{L}(a_i^\dagger)^{n_i}\right)\ket{0}
&&\text{w.\ }n_i=0,1,
\end{align}
using $[Z_i,Z_j]=0\ \forall i,j$, $[Z_i,c_j]=0$ for $i\neq j$ (hence
$[Z_i,a_j^\dagger]=0$ for $i\neq j$), and $Z_j\ket{0}=\ket{0}$. The ordering is not
important anymore because $[a_i^\dagger,a_j^\dagger]=0$.

\subsection{Action on basis states}

One cannot define the action of $a_j\,(a_j^\dagger)$ locally, since the JW string
$Z_1\cdots Z_{j-1}$ is non-local, so one has to examine them on the full many-body basis.
For $n_\ell=0$,
\begin{align}
a_\ell\ket{n_1,\dots,n_\ell=0,\dots,n_L}
&=a_\ell\left(\prod_{i=1}^{L}(a_i^\dagger)^{n_i}\right)\ket{0}\notag\\
&\overset{[a_\ell,a_j^\dagger]=0}{=}\left(\prod_{i=1}^{L}(a_i^\dagger)^{n_i}\right)a_\ell\ket{0}=0,
\end{align}
since $a_\ell\ket{0}=c_\ell\ket{0}=0$. For $n_\ell=1$,
\begin{align}
a_\ell\ket{n_1,\dots,n_\ell=1,\dots,n_L}
&=a_\ell\left(\prod_{\substack{i=1\\i\neq\ell}}^{L}(a_i^\dagger)^{n_i}\right)a_\ell^\dagger\ket{0}\notag\\
&\overset{[a_\ell,a_j^\dagger]=0}{=}
\left(\prod_{\substack{i=1\\i\neq\ell}}^{L}(a_i^\dagger)^{n_i}\right)a_\ell a_\ell^\dagger\ket{0}
=\left(\prod_{\substack{i=1\\i\neq\ell}}^{L}(a_i^\dagger)^{n_i}\right)c_\ell c_\ell^\dagger\ket{0}\notag\\
&=\left(\prod_{\substack{i=1\\i\neq\ell}}^{L}(a_i^\dagger)^{n_i}\right)\ket{0}
=\ket{n_1,\dots,n_\ell=0,\dots,n_L}.
\end{align}
Using $[a_\ell^\dagger,a_j^\dagger]=0$ and $(a_\ell^\dagger)^2=0$, one quickly finds
\begin{align}
a_\ell^\dagger\ket{n_1,\dots,n_\ell=0,\dots,n_L}&=\ket{n_1,\dots,n_\ell=1,\dots,n_L},\\
a_\ell^\dagger\ket{n_1,\dots,n_\ell=1,\dots,n_L}&=0.
\end{align}
Thus $a_\ell\,(a_\ell^\dagger)$ destroys (creates) a particle in the many-body state
\emph{without} giving rise to an extra sign structure (as $c_\ell\,(c_\ell^\dagger)$ would do).

\subsection{Identification with spins}

We can thus identify $n_i=0,1$ with $n_i=\,\downarrow,\uparrow$ in the many-body state, and the
new operators can be associated with spins by identifying
\begin{equation}
a_j=S_j^-,\qquad a_j^\dagger=S_j^+\qquad\&\qquad n_j-\tfrac12=S_j^z=\tfrac12\,\sigma_j^z,
\end{equation}
since $[S_j^-,S_j^+]=-2S_j^z$ for spins (see the quick refresher on spin,
\cref{sec:spin-refresher}). For instance
\begin{equation*}
a_3^\dagger\ket{0}=S_3^+\ket{\downarrow\downarrow\downarrow}=\ket{\uparrow\downarrow\downarrow}.
\end{equation*}
Matrix elements do not change under the JWT, since it is only a rewriting of operators and a
renaming of basis states.

\subsection{Proofs about the JWT} \label{subsec:proofs_jwt}

\paragraph{Claim: $\{Z_j,c_j^{(+)}\}=0$.}
\begin{proof}
On the single site $j$, with $Z_j\ket{0}_j=\ket{0}_j$, $Z_j\ket{1}_j=-\ket{1}_j$,
$c_j\ket{0}_j=0$ and $c_j\ket{1}_j=\ket{0}_j$,
\begin{align}
Z_j c_j\ket{0}_j&=Z_j\cdot 0=0=c_j\ket{0}_j=c_j Z_j\ket{0}_j,\\
Z_j c_j\ket{1}_j&=Z_j\ket{0}_j=\ket{0}_j=c_j\ket{1}_j,\\
c_j Z_j\ket{1}_j&=-c_j\ket{1}_j=-\ket{0}_j.
\end{align}
Hence, on both basis states $\{Z_j,c_j\}=Z_j c_j+c_j Z_j=c_j-c_j=0$, and since
$Z_j^\dagger=Z_j$ it follows that $\{Z_j,c_j^\dagger\}=0$ as well. (see Coleman, Chapter 4.1)
\end{proof}

\paragraph{Claim: $[a_i^{(+)},a_j^{(+)}]=0\ \, \forall\,i,j$.}
\begin{proof}
Assume w.l.o.g.\ that $i<j$. Using \cref{eq:jw-inv},
\begin{align}
a_j a_i&=\left(Z_1\cdots Z_{j-1}c_j\right)\left(Z_1\cdots Z_{i-1}c_i\right)\notag\\
&\overset{\substack{[Z_j,c_\ell]=0,\,j\neq\ell\\ \{c_i,c_j\}=0}}{=}
Z_1^2\cdots Z_{i-1}^2\,Z_i\cdots Z_{j-1}\,c_j c_i
=-Z_i\cdots Z_{j-1}\,c_i c_j\qquad(=0\text{ if }i=j),
\end{align}
and likewise
\begin{align}
a_i a_j&=\left(Z_1\cdots Z_{i-1}c_i\right)\left(Z_1\cdots Z_{j-1}c_j\right)\notag\\
&\overset{[Z_i,c_\ell]=0,\,i\neq\ell}{=}Z_{i+1}\cdots Z_{j-1}\,c_i Z_i c_j
\overset{\{Z_i,c_i\}=0}{=}-Z_i\cdots Z_{j-1}\,c_i c_j.
\end{align}
Comparing the two, $a_i a_j=a_j a_i$, i.e.\ $[a_i,a_j]=0$. Analogously for
$[a_i^\dagger,a_j^\dagger]$ and for $i>j$; the case $i=j$ is straightforward since
$c_i^2=0$.
\end{proof}

\paragraph{Claim: $[a_i,a_j^\dagger]=0\ \, \forall\,i\neq j$.}
\begin{proof}
We can show it w.l.o.g.\ for $i<j$. Using \cref{eq:jw-inv} and \cref{eq:jw-Zc},
\begin{align}
a_i^\dagger a_j&=\bigl(c_i^\dagger Z_{i-1}\cdots Z_1\bigr)\bigl(Z_1\cdots Z_{j-1}c_j\bigr)
\overset{Z_i^2=1,\,i<j}{=}c_i^\dagger\,Z_i\cdots Z_{j-1}\,c_j\notag\\
&=Z_{i+1}\cdots Z_{j-1}\,c_i^\dagger Z_i c_j
\overset{\{Z_i,c_i^\dagger\}=0}{=}-Z_i\cdots Z_{j-1}\,c_i^\dagger c_j
=+Z_i\cdots Z_{j-1}\,c_j c_i^\dagger,
\end{align}
and
\begin{align}
a_j a_i^\dagger &= \bigl(Z_1\cdots Z_{j-1}c_j\bigr)\bigl(c_i^\dagger Z_{i-1}\cdots Z_1\bigr) \notag \\
&= Z_i\cdots Z_{j-1}\,c_j c_i^\dagger\,Z_{i-1}^2\cdots Z_1^2 \qquad ([Z_j,c_\ell^{(+)}]=0,\,j\neq\ell) \notag \\
&=Z_i\cdots Z_{j-1}\,c_j c_i^\dagger .
\end{align}
Both reduce to the same expression, so $a_j a_i^\dagger-a_i^\dagger a_j=0$ for $i<j$, i.e.\
$[a_j,a_i^\dagger]=0$; the case $i>j$ follows straightforwardly.
\end{proof}

\paragraph{Claim: $[a_i,a_i^\dagger]=1-2n_i$ with $n_i=c_i^\dagger c_i$.}
\begin{proof}
Using \cref{eq:jw-inv} and $Z_k^2=\mathds 1$,
\begin{align}
a_i^\dagger a_i&=c_i^\dagger\,Z_{i-1}\cdots Z_1 Z_1\cdots Z_{i-1}\,c_i=c_i^\dagger c_i=n_i,\\
a_i a_i^\dagger&=Z_1\cdots Z_{i-1}\,c_i c_i^\dagger\,Z_{i-1}\cdots Z_1=c_i c_i^\dagger=1-c_i^\dagger c_i,
\end{align}
so that
\begin{equation}
[\hat a_i,\hat a_i^\dagger]=a_i a_i^\dagger-a_i^\dagger a_i=(1-c_i^\dagger c_i)-c_i^\dagger c_i
=1-2\hat n_i.
\end{equation}
\end{proof}

\end{document}
